\documentclass[a4paper,oneside,11pt]{book}

% ==============================================================================
% MASTER TEMPLATE LAPORAN PRAKTIKUM (VERSI MONOLITIK)
% ==============================================================================
% Dokumen ini memuat seluruh struktur laporan praktikum secara lengkap dan utuh.
% Ganti seluruh teks placeholder [DI DALAM KURUNG SIKU] dengan informasi riil.
%
% CATATAN MODULARITAS:
% Jika ingin memecah file ini menjadi bab-bab terpisah, Anda dapat memindahkan
% isi setiap bab ke file 'template-*.tex' masing-masing, lalu memanggilnya
% menggunakan perintah \input{template-namabab.tex}.
% ==============================================================================

% ==========================================
% 1. DEPENDENCIES DAN PAKET
% ==========================================
% Paket grafis dan manipulasi halaman
\usepackage{graphicx}
\usepackage{titling}
\usepackage{geometry}
\usepackage{float}
\usepackage{adjustbox}
\usepackage{indentfirst}
\usepackage{parskip}

% Paket tabel tingkat lanjut
\usepackage{array}
\usepackage{tabularx}
\usepackage{multirow}
\usepackage{booktabs}
\usepackage[table]{xcolor}

% Bahasa dan font
\usepackage[indonesian]{babel}
\usepackage[T1]{fontenc}
\usepackage{textcomp}
\usepackage{mathptmx}   
\usepackage{helvet}    
\usepackage{amsmath}
\usepackage{amsfonts}   

% Diagram TikZ (opsional untuk alur diagram)
\usepackage{tikz}
\usetikzlibrary{shapes.geometric, arrows}

% Header dan Footer kustom
\usepackage{fancyhdr}
\pagestyle{fancy}
\fancyhf{}
\rfoot{\thepage}
\renewcommand{\headrulewidth}{0pt}

% Pengaturan Margin Standar (3cm Kiri, 2.5cm Kanan, 3cm Atas, 3cm Bawah)
\geometry{
    left=3cm,
    right=2.5cm,
    top=3cm,
    bottom=3cm
}

% Hyperlink dan Referensi Silang
\usepackage{url}
\usepackage{hyperref}
\usepackage[all]{hypcap}
\hypersetup{
    colorlinks=true,
    linkcolor=black,
    citecolor=blue,
    urlcolor=blue
}

% Paket Kode Program (Listings)
\usepackage{listings}
\usepackage{etoolbox}

% Definisi Palet Warna Kode Program
\definecolor{codegreen}{rgb}{0,0.6,0}
\definecolor{codegray}{rgb}{0.5,0.5,0.5}
\definecolor{codepurple}{rgb}{0.58,0,0.82}
\definecolor{backcolour}{rgb}{0.95,0.95,0.92}

% Konfigurasi Gaya Kode Program Default
\lstdefinestyle{codestyle}{
    backgroundcolor=\color{backcolour},   
    commentstyle=\color{codegreen},
    keywordstyle=\color{codepurple},
    numberstyle=\tiny\color{codegray},
    stringstyle=\color{codegreen},
    basicstyle=\ttfamily\footnotesize,
    breakatwhitespace=false,         
    breaklines=true,                 
    captionpos=b,                    
    keepspaces=true,                 
    numbers=left,                    
    numbersep=5pt,                  
    showspaces=false,                
    showstringspaces=false,
    showtabs=false,                  
    tabsize=2,
    frame=single
}
\lstset{style=codestyle}

\begin{document}

% ==============================================================================
% 2. HALAMAN JUDUL (COVER)
% ==============================================================================
\begin{titlingpage}
\begin{center}
\vspace*{0.5cm}

% Judul Laporan
{\large \textbf{LAPORAN PRAKTIKUM}} \\[0.3cm]
{\large \textbf{[NAMA MATA KULIAH LENGKAP]}} \\[0.3cm]
{\large \textbf{PERTEMUAN [NOMOR PERTEMUAN]}} \\[0.3cm]
{\large \textbf{[JUDUL / TOPIK MATERI PRAKTIKUM]}} \\[2.5cm]

% Logo Universitas (Pastikan file lambangugm.png berada di folder yang sama)
\includegraphics[width=6cm]{lambangugm.png} \\[2.5cm]

% Identitas Penyusun
\normalsize
Disusun oleh:\\[0.5cm]
\begin{tabular}{@{}l@{\hspace{1cm}}l@{}}
Nama           & : [Nama Lengkap Mahasiswa] \\
NIM            & : [NIM Mahasiswa] \\
Kelas          & : [Kelas / Rombel] \\
Dosen Pengampu & : [Nama Dosen Pengampu, Gelar]
\end{tabular}

\vfill 

% Institusi
{\large \textbf{PROGRAM STUDI D-IV TEKNOLOGI REKAYASA PERANGKAT LUNAK}} \\[0.1cm]
{\large \textbf{DEPARTEMEN TEKNIK ELEKTRO DAN INFORMATIKA}} \\[0.1cm]
{\large \textbf{SEKOLAH VOKASI}} \\[0.1cm]
{\large \textbf{UNIVERSITAS GADJAH MADA}} \\[0.1cm]
{\large \textbf{YOGYAKARTA}} \\[0.1cm]
{\large \textbf{[TAHUN PENYUSUNAN, misal: 2026]}}

\vspace*{0.5cm}
\end{center}
\end{titlingpage}

% ==============================================================================
% 3. DAFTAR ISI & DAFTAR GAMBAR
% ==============================================================================
\tableofcontents
\listoffigures
\newpage

% ==============================================================================
% 4. BAB I: TUJUAN PRAKTIKUM
% ==============================================================================
\chapter{Tujuan Praktikum}
\label{chap:tujuanpraktikum}

Adapun tujuan dilaksanakannya praktikum pada modul ini adalah sebagai berikut:
\begin{enumerate}
    \item Mahasiswa mampu memahami konsep fundamental, arsitektur dasar, serta prinsip kerja dari [Topik / Teknologi Utama Praktikum].
    \item Mahasiswa mampu mengonfigurasi lingkungan pengembangan, menginisialisasi pustaka kerja, serta memvalidasi alur pemrosesan program secara benar.
    \item Mahasiswa dapat mengimplementasikan metode [Nama Metode / Algoritma] untuk memecahkan studi kasus atau penugasan praktikum yang diberikan.
    \item Mahasiswa mampu mengevaluasi luaran sistem, membedah karakteristik performa, serta menarik kesimpulan berbasis bukti empiris pengujian.
\end{enumerate}

% ==============================================================================
% 5. BAB II: DASAR TEORI
% ==============================================================================
\chapter{Dasar Teori}
\label{chap:dasarteori}

\section{Konsep Fundamental [Nama Topik 1]}
[Nama Topik 1] merupakan salah satu pilar utama dalam [Bidang Keilmuan] yang berfungsi untuk [uraikan fungsi dan urgensi konsep secara mendalam] [\ref{modulPraktikum}]. Melalui mekanisme ini, pengembang atau analis data dapat mengoptimalkan alur kerja sistem menjadi lebih terstruktur dan efisien [\ref{bukuTeks}]. Kelebihan utama pendekatan ini mencakup skalabilitas dan efisiensi alokasi memori, meskipun tetap memerlukan pertimbangan terkait kompleksitas komputasi yang terlibat [\ref{jurnalIlmiah}].

\begin{figure}[H]
    \centering
    \includegraphics[width=0.6\textwidth]{diagram_konsep.png}
    \caption{Diagram Arsitektur [Nama Topik 1] (Sumber: [\ref{modulPraktikum}])}
    \label{fig:diagram_konsep}
\end{figure}

\section{Metode dan Formulasi Matematis [Nama Metode 2]}
Metode [Nama Metode 2] dirumuskan untuk menghitung dan mengevaluasi keterkaitan antara variabel masukan dengan nilai luaran yang diharapkan. Persamaan matematis dari metode ini dapat dinyatakan sebagai berikut [\ref{jurnalIlmiah}]:

\begin{equation}
    y = f(X) = \sum_{i=1}^{n} w_i \cdot x_i + b
    \label{eq:formulasi_dasar}
\end{equation}

Dimana:
\begin{itemize}
    \item $y$ = nilai luaran prediksi / estimasi sistem
    \item $x_i$ = nilai variabel masukan atau fitur ke-$i$
    \item $w_i$ = koefisien bobot pengaruh (\textit{weight})
    \item $b$ = konstanta bias (\textit{intercept})
    \item $n$ = jumlah total variabel prediktor
\end{itemize}

\section{Metrik Evaluasi Kinerja}
Guna mengevaluasi ketepatan luaran eksperimen secara objektif, digunakan metrik kesalahan galat numerik yang dirumuskan melalui persamaan:

\begin{equation}
    E = \sqrt{\frac{1}{m}\sum_{j=1}^{m}\left(y_{j} - \hat{y}_{j}\right)^2}
    \label{eq:metrik_evaluasi}
\end{equation}

Tingkat kesalahan yang semakin kecil membuktikan bahwa hasil pemodelan atau eksekusi program memiliki akurasi yang tinggi serta mampu mengikuti dinamika data aktual secara konsisten [\ref{jurnalIlmiah}].

% ==============================================================================
% 6. BAB III: HASIL DAN PEMBAHASAN
% ==============================================================================
\chapter{Hasil dan Pembahasan}
\label{chap:hasilpembahasan}

% ==========================================
% SEKSI 1: HASIL PRAKTIKUM
% ==========================================
\section{Hasil Praktikum}
\label{sec:hasilpraktikum}
Pada bagian Hasil Praktikum, dilakukan serangkaian eksperimen praktikum terbimbing sesuai dengan panduan modul [Nama Topik Praktikum], mulai dari penyiapan dependensi hingga pengujian luaran fungsi utama.

\subsection{Inisialisasi Lingkungan dan Import Library}
Tahap awal dimulai dengan memuat seluruh modul dan pustaka yang dibutuhkan selama sesi praktikum ke dalam lingkungan kerja.

\begin{figure}[H]
    \centering
    \includegraphics[width=0.55\textwidth]{1.png}
    \caption{Source Code Import Library dan Inisialisasi Lingkungan}
    \label{fig:1}
\end{figure}
Penjelasan Gambar :
\begin{itemize}
    \item \textbf{Input:} Mengeksekusi perintah import pustaka pemrosesan data dan visualisasi grafik yang dibutuhkan oleh modul praktikum.
    \item \textbf{Output:} Seluruh pustaka berhasil dimuat ke dalam memori lingkungan kerja tanpa terjadi kendala dependensi atau galat sistem.
\end{itemize}

\subsection{Prapemrosesan Data dan Transformasi Fitur}
Data masukan diproses dan disesuaikan formatnya agar kompatibel dengan modul komputasi utama.

\begin{figure}[H]
    \centering
    \includegraphics[width=0.50\textwidth]{2a.png}
    \caption{Source Code Prapemrosesan Data}
    \label{fig:2a}
    
    \vspace{0.2cm}
    
    \includegraphics[width=0.50\textwidth]{2b.png}
    \caption{Output Verifikasi Struktur Data}
    \label{fig:2b}
\end{figure}
Penjelasan Gambar :
\begin{itemize}
    \item \textbf{Input:} Pada Gambar \ref{fig:2a}, baris program mengeksekusi tahapan pembersihan dan penyesuaian dimensi variabel data masukan.
    \item \textbf{Output:} Gambar \ref{fig:2b} memvalidasi bahwa bentuk matriks dan tipe data telah sesuai dengan spesifikasi fungsi pengolahan inti.
\end{itemize}

\subsection{Visualisasi dan Evaluasi Luaran Terbimbing}
Luaran dari implementasi terbimbing dianalisis secara grafis untuk memahami pola respons sistem.

\begin{figure}[H]
    \centering
    \begin{minipage}{0.48\textwidth}
        \centering
        \includegraphics[width=\textwidth]{3a.png}
        \caption{Visualisasi Luaran Metode 1}
        \label{fig:3a}
    \end{minipage}\hfill
    \begin{minipage}{0.48\textwidth}
        \centering
        \includegraphics[width=\textwidth]{3b.png}
        \caption{Visualisasi Luaran Metode 2}
        \label{fig:3b}
    \end{minipage}
\end{figure}
Penjelasan Gambar :
\begin{itemize}
    \item \textbf{Output:} Gambar \ref{fig:3a} dan Gambar \ref{fig:3b} menampilkan komparasi visual luaran antar metode secara berdampingan, di mana Metode 2 memperlihatkan respons yang lebih stabil dalam menangani fluktuasi data.
\end{itemize}

% ==========================================
% SEKSI 2: TUGAS PRAKTIKUM
% ==========================================
\clearpage
\section{Tugas Praktikum}
\label{sec:tugaspraktikum}
Pada bagian Tugas Praktikum, diselesaikan studi kasus mandiri untuk menguji pemahaman mahasiswa terhadap materi yang telah diajarkan menggunakan dataset atau skenario baru.

\subsection{Implementasi Solusi Tugas Mandiri}
Penerapan konsep dilakukan pada dataset penugasan dengan konfigurasi parameter yang disesuaikan.

\begin{figure}[H]
    \centering
    \includegraphics[width=0.55\textwidth]{T1.png}
    \caption{Source Code Implementasi Solusi Tugas Praktikum}
    \label{fig:t1}
\end{figure}
Penjelasan Gambar :
\begin{itemize}
    \item \textbf{Input:} Mengeksekusi modul solusi yang dirancang secara mandiri untuk menyelesaikan permasalahan pada studi kasus penugasan.
    \item \textbf{Output:} Program berhasil dieksekusi secara lengkap dan menghasilkan luaran komputasi yang konsisten terhadap kriteria penugasan.
\end{itemize}

\subsection{Evaluasi Performa Tugas Mandiri}
Hasil pengujian dievaluasi untuk mengukur performa solusi secara kuantitatif.

\begin{figure}[H]
    \centering
    \includegraphics[width=0.60\textwidth]{T2.png}
    \caption{Output Hasil Pengujian dan Evaluasi Tugas Praktikum}
    \label{fig:t2}
\end{figure}
Penjelasan Gambar :
\begin{itemize}
    \item \textbf{Input:} Perintah penghitungan metrik performa terhadap hasil keluaran tugas praktikum.
    \item \textbf{Output:} Menampilkan ringkasan metrik evaluasi yang membuktikan bahwa solusi yang diterapkan berhasil mencapai target presisi yang diharapkan.
\end{itemize}

\subsection{Pembahasan dan Analisis Komparatif}
Berdasarkan hasil pengerjaan tugas mandiri, diperoleh temuan analitis sebagai berikut:
\begin{enumerate}
    \item \textbf{Pengaruh Konfigurasi Parameter:} Pemilihan parameter yang tepat terbukti meningkatkan ketahanan sistem terhadap derau data.
    \item \textbf{Efisiensi Solusi:} Algoritma yang diusulkan mampu menyelesaikan komputasi dengan efisiensi waktu yang dapat diterima.
\end{enumerate}

% ==============================================================================
% 7. BAB IV: KESIMPULAN
% ==============================================================================
\chapter{Kesimpulan}
\label{chap:kesimpulan}

Berdasarkan serangkaian kegiatan praktikum mengenai [Topik Praktikum] yang telah dilaksanakan—meliputi pemodelan langkah terbimbing pada Hasil Praktikum hingga penyelesaian studi kasus mandiri pada Tugas Praktikum—diperoleh beberapa simpulan penting sebagai berikut:

\begin{enumerate}
    \item \textbf{Efektivitas Implementasi Metode Utama:} Praktikum ini membuktikan bahwa metode [Nama Metode / Teknologi] mampu mengolah dan menyelesaikan alur kerja sistem secara terstruktur dan terukur sesuai dengan target capaian modul.
    
    \item \textbf{Pentingnya Tahapan Prapemrosesan Data:} Ketepatan penyiapan parameter lingkungan dan prapemrosesan data terbukti krusial dalam meminimalisasi galat eksekusi dan meningkatkan kualitas luaran akhir.
    
    \item \textbf{Keunggulan Hasil Pengujian:} Berdasarkan hasil pengujian dan visualisasi, pendekatan yang diimplementasikan mampu memberikan luaran yang stabil dan akurat sesuai indikator metrik evaluasi.
    
    \item \textbf{Penyelesaian Studi Kasus Mandiri:} Eksperimen pada tugas praktikum mengonfirmasi bahwa solusi yang dirancang mampu beradaptasi dengan baik pada skenario data baru tanpa mengalami penurunan performa signifikan.
    
    \item \textbf{Pengembangan Kompetensi Praktis:} Pemahaman terhadap alur kerja praktikum ini memberikan bekal pengetahuan aplikatif yang relevan untuk diterapkan pada perancangan sistem komputasi di masa mendatang.
\end{enumerate}

% ==============================================================================
% 8. BAB V: DAFTAR PUSTAKA
% ==============================================================================
\chapter{Daftar Pustaka}
\label{chap:daftarpustaka}

\begin{enumerate}
    \item \label{modulPraktikum} [Nama Dosen Pengampu, Gelar], dan Tim Asisten. \textit{Modul Praktikum [Nama Mata Kuliah]: [Topik Praktikum Pertemuan X]}. Yogyakarta: Departemen Teknik Elektro dan Informatika, Sekolah Vokasi, Universitas Gadjah Mada, [Tahun].
    
    \item \label{bukuTeks} [Nama Penulis Buku, Inisial]. \textit{[Judul Buku Referensi: Subjudul]}, [Edisi] ed. [Kota Penerbit]: [Nama Penerbit], [Tahun].
    
    \item \label{jurnalIlmiah} [Penulis 1] and [Penulis 2], "[Judul Artikel Ilmiah / Riset Terkait]," \textit{[Nama Jurnal Ilmiah]}, vol. [Volume], no. [Nomor], pp. [Awal]--[Akhir], [Tahun]. [Online]. Tersedia: \url{https://doi.org/...}.
    
    \item \label{sumberDataset} [Penyedia / Penulis Dataset], \textit{[Nama Dataset / Kumpulan Data]}. Kaggle / Repositori Data, [Tahun]. [Online]. Tersedia: \url{https://...}.
    
    \item \label{sourceCode} [Nama Mahasiswa], \textit{[Kode Matkul]-[NIM] (Source Code Laporan Praktikum)}. Google Colaboratory / GitHub, [Tahun]. [Online]. Tersedia: \url{https://colab.research.google.com/drive/...}.
\end{enumerate}

\end{document}
